\chapter{Python essentials for health professionals}
\label{ch:python-essentials}

\begin{quote}
\emph{``The goal is not to learn Python. The goal is to answer health questions with data. Python is just the tool.''}
\end{quote}

% ============================================================
\section{Why Python for health data?}

Three reasons health professionals should learn Python rather than relying solely on Excel or SPSS:

\begin{enumerate}
  \item \textbf{Reproducibility.} A Python script documents every step of your analysis. When a reviewer asks ``how did you handle missing lab values?'', you show the code --- not a description of menu clicks.
  \item \textbf{Scale.} Excel struggles with 100\,000 patient records. Python handles millions without slowing down.
  \item \textbf{Ecosystem.} Libraries like \texttt{lifelines} (survival analysis), \texttt{geopandas} (disease mapping), and \texttt{scikit-learn} (predictive models) give you tools that SPSS does not offer.
\end{enumerate}

\begin{clinicalnote}
You will not need to memorize syntax. Health professionals use Python like clinicians use a stethoscope --- you learn the parts you need, and you look up the rest. Every code block in this course can be copied into a Jupyter notebook and run immediately.
\end{clinicalnote}

% ============================================================
\section{Setting up your environment}

\subsection{Google Colab (recommended)}

No installation. Go to \url{https://colab.research.google.com}, sign in with a Google account, and create a new notebook. Every library we use is pre-installed.

\subsection{Local installation with Anaconda}

If you prefer working offline:
\begin{enumerate}
  \item Download Anaconda: \url{https://www.anaconda.com/download}
  \item Install with default settings.
  \item Open \textbf{Jupyter Notebook} from Anaconda Navigator.
\end{enumerate}

% ============================================================
\section{Your first notebook}

A Jupyter notebook is a document that mixes text, code, and output. Each \emph{cell} contains either:
\begin{itemize}
  \item \textbf{Code} --- Python instructions that produce a result.
  \item \textbf{Markdown} --- formatted text (explanations, headings, notes).
\end{itemize}

Press \texttt{Shift + Enter} to run a cell.

\begin{minted}{python}
# Your first Python cell
print("Hello, health data!")
\end{minted}

% ============================================================
\section{Variables and types}

A \emph{variable} stores a value. In health data, you work with four types constantly:

\begin{minted}{python}
patient_age = 45              # int (integer)
temperature = 37.2            # float (decimal number)
diagnosis = "Type 2 Diabetes" # str (string / text)
is_hospitalized = True        # bool (True or False)
\end{minted}

\begin{pythontip}
Variable names should be descriptive. Use \texttt{patient\_age}, not \texttt{x}. Your future self (and your collaborators) will thank you.
\end{pythontip}

% ============================================================
\section{Collections: lists and dictionaries}

\subsection{Lists}

A list holds multiple values in order:
\begin{minted}{python}
blood_pressures = [120, 135, 128, 142, 118]
print(f"Number of readings: {len(blood_pressures)}")
print(f"First reading: {blood_pressures[0]}")  # indexing starts at 0
print(f"Last reading: {blood_pressures[-1]}")
\end{minted}

\subsection{Dictionaries}

A dictionary maps keys to values --- like a patient chart:
\begin{minted}{python}
patient = {
    "id": "PAT-0042",
    "age": 58,
    "sex": "F",
    "diagnosis": "Hypertension",
    "systolic_bp": 148,
    "medications": ["Amlodipine", "Hydrochlorothiazide"]
}
print(patient["diagnosis"])  # Hypertension
\end{minted}

% ============================================================
\section{Control flow}

\subsection{Conditionals}

\begin{minted}{python}
systolic = 148

if systolic >= 140:
    category = "Stage 2 Hypertension"
elif systolic >= 130:
    category = "Stage 1 Hypertension"
elif systolic >= 120:
    category = "Elevated"
else:
    category = "Normal"

print(f"BP category: {category}")
\end{minted}

\subsection{Loops}

\begin{minted}{python}
patients = [
    {"id": "P001", "age": 34, "hba1c": 5.4},
    {"id": "P002", "age": 67, "hba1c": 7.8},
    {"id": "P003", "age": 52, "hba1c": 6.1},
]

for p in patients:
    if p["hba1c"] >= 6.5:
        print(f"{p['id']}: Diabetic (HbA1c = {p['hba1c']})")
    elif p["hba1c"] >= 5.7:
        print(f"{p['id']}: Pre-diabetic (HbA1c = {p['hba1c']})")
    else:
        print(f"{p['id']}: Normal (HbA1c = {p['hba1c']})")
\end{minted}

% ============================================================
\section{Functions}

A function encapsulates a reusable computation:

\begin{minted}{python}
def bmi(weight_kg, height_m):
    """Calculate Body Mass Index."""
    return weight_kg / (height_m ** 2)

def bmi_category(bmi_value):
    """Classify BMI according to WHO categories."""
    if bmi_value < 18.5:
        return "Underweight"
    elif bmi_value < 25.0:
        return "Normal weight"
    elif bmi_value < 30.0:
        return "Overweight"
    else:
        return "Obese"

value = bmi(82, 1.75)
print(f"BMI: {value:.1f} ({bmi_category(value)})")
\end{minted}

\begin{exercise}
\begin{enumerate}
  \item Write a function \texttt{classify\_bp(systolic, diastolic)} that returns a blood pressure category according to the AHA/ACC guidelines.
  \item Create a list of 5 patient dictionaries with \texttt{weight\_kg} and \texttt{height\_m} fields. Loop through them, compute BMI, and print the category for each.
  \item Write a function \texttt{egfr(creatinine, age, sex)} that estimates glomerular filtration rate using the CKD-EPI equation.
\end{enumerate}
\end{exercise}

% ============================================================
\section{Importing libraries}

Python's power comes from libraries. Here are the ones we will use throughout this course:

\begin{minted}{python}
import pandas as pd          # data manipulation
import numpy as np           # numerical computing
import matplotlib.pyplot as plt  # plotting
import seaborn as sns        # statistical visualization
from scipy import stats      # statistical tests
\end{minted}

\begin{warningbox}
If you get \texttt{ModuleNotFoundError}, install the missing library:
\begin{minted}{bash}
pip install pandas matplotlib seaborn scipy
\end{minted}
In Google Colab, prefix with \texttt{!}: \texttt{!pip install lifelines}
\end{warningbox}

% ============================================================
\section{Chapter summary}

\begin{itemize}
  \item Python is free, reproducible, and has health-specific libraries.
  \item Jupyter notebooks mix code, text, and output in one document.
  \item Variables store values; lists and dictionaries organize them.
  \item Functions make your analysis reusable and readable.
  \item \texttt{pandas}, \texttt{matplotlib}, \texttt{seaborn}, and \texttt{scipy} are the core libraries for health data analysis.
\end{itemize}
